\chapter{Outils logiciels}

\mysetref{tools}%
\begin{enumerate}

\item{Cellsim}\label{soft-cellsim}
  - CA Simulator\hfil\break
  Cellsim copyright 1989, 1990 by Chris Langton and Dave Hiebeler
  (\email{cgl}{lanl.gov}, \email{hiebeler}{heretic.lanl.gov})\hfil\break
  Version 2.5 of Cellsim is a SunView-based cellular automata simulator
  allowing interactive specification, editing, running, and analysis of
  1- and 2-D CA's.  It will run on Sun-3's, -4's, and Sparc stations,
  color or B\&W.

\item{Scamper}\label{soft-scamper}
  - a simulator of Cellular automata.\hfil\break
  Copyright 1991 by Ian Palmer of Imperial College of Science, Technology
  and Medicine, University of London.

\item{CA LAB}\label{soft-ca_lab}
  \hfil\break
  Copyright 1989 Rudy Rucker and Autodesk. Inc.

\item{Xlife}\label{soft-xlife}
  - Conway's Game of Life for X, version 3.0\hfil\break XLife
  Copyright 1989 Jon Bennett \\
  \email{jb7m+}{andrew.cmu.edu}, \email{jcrb}{cs.cmu.edu}

\item{HODGE-C}\label{soft-hodge_c}
  - A C implementation of Gerhard \& Schuster's hodge-podge machine\hfil\break
  Copyright (C) 1993 Joerg Heitkoetter

\item{Cellang}\label{soft-cellang}
  \hfil\break
  Copyright (C) 1992  J Dana Eckart\hfil\break
  cellc - Cellang 2.0 (cellular automata) compiler\hfil\break
  cellview - cellular automata viewer\hfil\break
  pe-scam - Cellang 2.0 (cellular automata) machine driver

\item{pbmplus}\label{soft-pbm}
  - Extended Portable Bitmap Toolkit\hfil\break
  Copyright (C) 1989, 1991 by Jef Poskanzer.

\item{groff}\label{soft-groff}
  - document formatting system\hfil\break
  J. Clark. Copyright (C) 1989, 1990, 1991, 1992 Free Software Foundation,
  Inc.\hfil\break
  From the README:\hfil\break
  Included in this release are implementations of troff, pic, eqn,
  tbl, refer, the -man macros and the -ms macros, and drivers for
  PostScript, TeX dvi format, and typewriter-like devices.  Also
  included is a modified version of the Berkeley -me macros, an
  enhanced version of the X11 xditview previewer, and an
  implementation of the -mm macros contributed by Joergen Haegg.

\item{ghostscript}\label{soft-ghostscript}
  - PostScript previewer\hfil\break
  Copyright (C) 1989, 1992 Aladdin Enterprises.  Distributed by Free
  Software Foundation, Inc.

\item{xv}\label{soft-xv}
  - interactive image display for the X Window System\hfil\break
  Copyright 1989, 1990, 1991, 1992 by John Bradley and the University
  of Pennsylvania

\item{gnuplot}\label{soft-gnuplot}
  - an interactive plotting program\hfil\break
  Copyright(C) 1986, 1987, 1990, 1991, 1992 Thomas Williams, Colin
  Kelley

\item{Khoros}\label{soft-Khoros}
  - an integrated software development environment for
  information processing and visualization.\hfil\break
  Khoros Consortium. Khoros is a Registered Trademark of The University of New
  Mexico\hfil\break
  From the README:\hfil\break
  Khoros components include a visual programming language, code
  generators for extending the visual language and adding new
  application packages to the system, an interactive user interface
  editor, interactive image display programs, surface visualization,
  an extensive library of image processing, numerical analysis and
  signal processing routines, and 2D/3D plotting packages.  X
  applications built using the Khoros User Interface Libraries have
  built in journal/playback and groupware capabilities.

\item{cantata}\label{soft-cantata}
  - visual programming environment for accessing the
    Khoros library of routines

\item{qrt}\label{soft-qrt}
  - Ray Tracer\hfil\break
  Copyright 1988 and 1989 Steve Koren

\item{vart}\label{soft-vart}
  - 3D object toolbox for Khoros\hfil\break
  Copyright 1992, Peter Averkamp.  All rights reserved.

\item{pgmtexture}\label{soft-pgmtexture}
  - calculate textural features on a portable graymap\hfil\break
  J. D. McCauley. Copyright (C) 1991 by Texas Agricultural Experiment
  Station

\item{gnans}\label{soft-gnans}
  - A program for Stochastic and Deterministic Dynamical Systems\hfil\break
  B. Martensson. Copyright (C) 1989, 1991 Free Software Foundation,
  Inc.

\item{Tcl}\label{soft-Tcl}
  - tool command language facilities\hfil\break
  J. Ousterhout. Copyright 1987-1991 Regents of the University of
  California

\item{Tk}\label{soft-Tk}
  - an X11 toolkit that provides the Motif look and feel and is
  implemented using the Tcl command language.\hfil\break
  J. Ousterhout. Copyright 1987-1991 Regents of the University of
  California

\item{XDataSlice}\label{soft-XDS}
  \hfil\break
  National Center for Supercomputing Applications, September 1989\\ %
  NCSA X DataSlice Version 1.0 source code and documentation are in
  the public domain.\\ %
  From documentation:\\ %
  NCSA X DataSlice is a color imaging and data analysis tool based on
  the X Window System. It was developed for analysis of
  three-dimensional 32-bit floating-point scientific data stored in
  NCSA HDF file format.

\item{IsoVis}\label{soft-isovis}
  \hfil\break
  National Center for Supercomputing Applications, July 1990\\ %
  NCSA Isosurface Visualizer Version 1.0 source code and documentation
  are in the public domain.\\ %
  From documentation:\\ %
  Isosurface Visualizer, a 3D visualization tool, allows users to view
  data in a simple manner. Input data consists of a single 3D scalar
  array of floating-point values. The indices of this array form the
  x, y, z coordinates, and the variable at each node determines where
  the isosurface passes through the data. You can think of an
  isosurface as a three-dimensional contour plot, although this is not
  what it actually is.

\item{HDF}\label{soft-HDF}
  \hfil\break
  National Center for Supercomputing Applications, November 8, 1993\\ %
  NCSA HDF Version 3.3 source code and documentation are in the public
  domain.\\ %
  From documentation:\\ %
  The Hierarchical Data Format (HDF) was designed to be an easy,
  straight-forward, and self-describing means of sharing scientific
  data among people, projects, and types of computers. An extensible
  header and carefully crafted internal layers provide a system that
  can grow as scientific data-handling needs evolve.

\item{VIS-5D VERSION 3.3}\label{soft-VIS-5D}
  \hfil\break
  Space Science and Engineering Center, University of Wisconsin -
  Madison, 1225 West Dayton Street, Madison, WI  53706\\ %
  From documentation:\\ %
  VIS-5D is a software system for visualizing data made by numerical
  weather models and similar sources.  VIS-5D works on data in the
  form of a five-dimensional rectangle.  That is, the data are real
  numbers at each point of a "grid" or "lattice" which spans three
  space dimensions, one time dimension and a dimension for enumerating
  multiple physical variables.

\item{viewit v3.13}\label{soft-viewit}
  \hfil\break
  NMR Imaging and Spectroscopy, April 20, 1993\\ %
  C.D.Gregory Biomedical Magnetic Resonance Laboratory College of
  Medicine University of Illinois at Urbana-Champaign Urbana Illinois\\ %
  From the README:\\ %
  Perhaps, the most useful description of viewit is a "scientific
  calculator" for multidimensional arrays with extra "buttons" for
  signal processing, scientific visualization, image processing and
  multivariate analysis.
  
\item{LIC}\label{soft-LIC}
  - Line Integral Convolution\hfil\break
  Brian Cabral and Casey Leedom. Copyright (c) 1993 The Regents of the
  University of California. All rights reserved.

\item{Octave}\label{soft-Octave}
  - a high-level language for numerical computations.\hfil\break
  Copyright (C) 1992, 1993, 1994 John W. Eaton

\item{RLaB}\label{soft-Rlab}
  - matrix oriented, interactive programming environment.\hfil\break
   Copyright (C) 1992  Ian R. Searle

\end{enumerate}

%%
%% Local Variables:
%% mode: auto-fill
%% iso-accents-minor-mode: nil
%% TeX-master: "top"
%% TeX-command-default: "mytex"
%% Local IspellParsing: "latex-mode ~latin1"
%% Local IspellPersDict: "~/.ispell_lisp"
%% End:
